\section{Revised latent-class analysis}
\label{sec:revised-lca}

The expanded revised comparison tested one through six classes in two model
families on the initial 407-case indicator set. Neither family yielded an
admissible multiclass fit. The later discrete-item sensitivity analysis below
tests a broader set of measurement choices. Observed profiles remain the
paper's basis for reporting.

\paragraph{Indicators and common sample.}
We used five reported items: regular use of a personal smartphone; proof of
address in one's own name that the respondent considered acceptable to a
financial institution; confidence in avoiding online fraud and suspicious
messages; deciding how household remittances are used most of the time; and
perceived information about fees, exchange rates and what to do if something
goes wrong. The first two items were binary. For each agreement item, we
combined adjacent degrees of disagreement and of agreement, retaining neutral
as its own category. The information item retained all four ordered responses.
The recodes apply only to this analysis. Of the 423 survey records, 407 had
scorable responses on all five items. The 11 respondents whose address
documentation was in process and the five who declined to answer that
question were excluded from this complete-case comparison, not coded as
lacking proof of address.

\paragraph{Specifications and checks.}
The local-independence model used two binary logits and three ordered logits
with class-specific cutpoints, treating responses as independent within class.
The local-dependence sensitivity modeled the nine-category joint response
to the two agreement items with a multinomial logit, allowing those items
to remain associated within class. The other items were unchanged. Both
specifications assign probabilities to the same 144 possible response
patterns on the same 407 records; 63 patterns occurred in the sample.
We computed a one-class baseline from sample frequencies for each family.
For two through six classes, we used three separately seeded fits at each
class count in each family, with 120 initialization draws and a final
iteration cap of 2,400 per fit in Stata/MP~19. The number of free parameters
is $10K-1$ under local independence and $14K-1$ under local dependence.
We fixed the search and reporting rules before the expanded run. This
was a retrospective sensitivity analysis, not a preregistered test.

Reporting a class solution required an interior, stationary fit with
adequate class sizes and a usable covariance matrix. At least two
independently seeded fits had to agree on the likelihood and profile
probabilities after accounting for arbitrary class labels. In
Table~\ref{tab:revised-lca}, the only four fits that Stata reported as
converged had two classes. All four had at least one probability within
$10^{-6}$ of zero or one. Two also failed the stationarity check: the
maximum discrepancy between posterior-weighted counts and their fitted
counterparts, divided by the sample size, exceeded $10^{-6}$. The other
26 fits reached the iteration cap. We reconstructed all 30 terminal
likelihoods independently from the saved coefficients and observed
pattern counts. Numerical agreement verifies the recorded endpoints;
it does not establish convergence.

\begin{table}[!htbp]
\centering
\caption{Class-count diagnostics for the revised latent-class models}
\label{tab:revised-lca}
\begin{minipage}{\linewidth}
\centering\small
\setlength{\tabcolsep}{4.5pt}
\begin{tabular}{@{}lrrrrrrr@{}}
\toprule
Model & $K$ & \shortstack{Free\\parameters} & \shortstack{Stata\\converged}
& \shortstack{At iteration\\cap} & \shortstack{Log\\likelihood} & AIC & BIC \\
\midrule
Local independence & 1 & 9 & --- & --- & -1398.41 & 2814.83 & 2850.91 \\
                   & 2 & 19 & 2 & 1 & -1372.17 & 2782.35 & 2858.51 \\
                   & 3 & 29 & 0 & 3 & -1357.59 & 2773.18 & 2889.43 \\
                   & 4 & 39 & 0 & 3 & -1350.46 & 2778.92 & 2935.26 \\
                   & 5 & 49 & 0 & 3 & -1342.92 & 2783.83 & 2980.26 \\
                   & 6 & 59 & 0 & 3 & -1339.11 & 2796.23 & 3032.75 \\
\midrule
Local dependence & 1 & 13 & --- & --- & -1392.14 & 2810.29 & 2862.40 \\
                   & 2 & 27 & 2 & 1 & -1363.78 & 2781.55 & 2889.79 \\
                   & 3 & 41 & 0 & 3 & -1352.42 & 2786.83 & 2951.19 \\
                   & 4 & 55 & 0 & 3 & -1344.60 & 2799.20 & 3019.69 \\
                   & 5 & 69 & 0 & 3 & -1338.52 & 2815.04 & 3091.65 \\
                   & 6 & 83 & 0 & 3 & -1334.24 & 2834.48 & 3167.21 \\
\bottomrule
\end{tabular}
\par\smallskip\footnotesize\raggedright
Notes: All rows use 407 complete cases. $K$ is the number of classes.
The one-class baselines use no optimizer. Other rows summarize three
fits each. The likelihood and the Akaike and Bayesian information criteria
(AIC and BIC) use the largest recorded endpoint in each row, not a verified
multiclass maximum. None of the multiclass fits passed the reporting checks.\par
Source: Authors' calculations from the CFI-DPI survey in Colombia.
\end{minipage}
\end{table}

\paragraph{Stability of the four-class candidates.}
None of the six four-class fits converged. For each pair of seeds we matched
class labels to minimize the largest difference across class shares and
response probabilities. Even the closest pair differed by about
20 percentage points on at least one quantity under local independence
and 28 points under local dependence. These are discrepancies between
terminal solutions, not intervals around a validated four-class model.
The historical fit in Section~\ref{sec:lca} uses a different set of
indicators and cannot settle the class count here.

\paragraph{Implications for reporting.}
BIC was lowest at the analytic one-class baseline within each family among
the recorded endpoints. AIC favored different endpoints in the two
families, but neither was admissible. There is also a bound that does
not depend on optimization: no model on these data can exceed the
saturated multinomial log likelihood, $\sum_{j:n_j>0} n_j\log(n_j/N)$,
where $n_j$ is the count for response pattern $j$. Using that bound in
BIC, even a four-class local-dependence fit attaining the saturated
likelihood would have a BIC of at least 2929.54, compared with 2862.40
for its one-class baseline. The bound also rules out a BIC improvement
at five or six classes in either family. It does not settle the smaller
candidates or rule out latent heterogeneity. No multiclass solution from this
comparison passed the reporting checks. The paper's four observed
onboarding profiles use a different pair of questions and 404 scorable
pairs; they are not estimates of latent groups.

\input{"\CFIPaperDir supplement/discrete-lca-sensitivity.tex"}
